\section{Punktweiser Vergleich mehrfach indizierter Mengen}

\FormulaRefAuto{B27TermImagePointwiseEquality}[thm] besagt, dass
punktweise gleiche Terme dieselbe Menge bilden. Für drei getrennte
Indexmengen wird dieses Prinzip nun unmittelbar mit den beschränkten
Existenzquantoren formuliert. Dadurch muss ein späterer Beweis die drei
Zeugen nicht nochmals einzeln wählen und wieder zusammensetzen.

Dabei sind \(t(x,y,z)\) und \(r(x,y,z)\) Mengenterme mit festen
Parametern. Die verschiedenen Bindungsvariablen \(x,y,z,w\) kommen
nicht in diesen Parametern oder in \(X,Y,Z,M,N\) frei vor; \(w\)
kommt auch in den beiden Termen nicht frei vor. Die Indexmengen dürfen
leer sein. Es werden weder Funktionen noch algebraische Operationen
vorausgesetzt.

\FormulaThmDeltaK[Punktweise gleiche Terme haben dieselben Dreifachzeugen]{%
  \begin{aligned}[t]
    &\forall x\in X\,\forall y\in Y\,\forall z\in Z\;
      t(x,y,z)=r(x,y,z)\vdash{}\\[-2pt]
    &\bigl(\exists x\in X\,\exists y\in Y\,\exists z\in Z\;
      w=t(x,y,z)\bigr)\\[-2pt]
    &\qquad\leftrightarrow
      \bigl(\exists x\in X\,\exists y\in Y\,\exists z\in Z\;
      w=r(x,y,z)\bigr)
  \end{aligned}
}{TripleTermWitnessPointwiseEquality}{%
  \DeltaRow{Mengen}{X\dsep Y\dsep Z\dsep x\dsep y\dsep z\dsep w}
  \DeltaRow{Mengenterme}{t(x,y,z)\dsep r(x,y,z)}
}
\begin{tabproofsplitwide}
  \proofpartwide{Vom ersten Term zum zweiten}
  \proofstepwidestar[1]{\forall x\in X\,\forall y\in Y\,
    \forall z\in Z\;t(x,y,z)=r(x,y,z)}{\rA}
  \proofstepwidestar[2]{\exists x\in X\,\exists y\in Y\,
    \exists z\in Z\;w=t(x,y,z)}{\rA}
  \proofstepwidestar[3]{x\in X\land\exists y\in Y\,
    \exists z\in Z\;w=t(x,y,z)}{\rA}
  \proofstepwidestar[3]{x\in X}{\rAEa{3}}
  \proofstepwidestar[3]{\exists y\in Y\,\exists z\in Z\;
    w=t(x,y,z)}{\rAEb{3}}
  \proofstepwidestar[6]{y\in Y\land\exists z\in Z\;
    w=t(x,y,z)}{\rA}
  \proofstepwidestar[6]{y\in Y}{\rAEa{6}}
  \proofstepwidestar[6]{\exists z\in Z\;w=t(x,y,z)}{\rAEb{6}}
  \proofstepwidestar[9]{z\in Z\land w=t(x,y,z)}{\rA}
  \proofstepwidestar[9]{z\in Z}{\rAEa{9}}
  \proofstepwidestar[9]{w=t(x,y,z)}{\rAEb{9}}
  \proofstepwidestar[1,3,6,9]{t(x,y,z)=r(x,y,z)}{%
    \rRE{10,\rUE{\rRE{7,\rUE{\rRE{4,\rUE{1}}}}}}}
  \proofstepwidestar[1,3,6,9]{w=r(x,y,z)}{\rIE{12,11}}
  \proofstepwidestar[1,3,6,9]{\exists x\in X\,\exists y\in Y\,
    \exists z\in Z\;w=r(x,y,z)}{%
    \rEI{\rAI{4,\rEI{\rAI{7,\rEI{\rAI{10,13}}}}}}}
  \proofstepwidestar[1,3,6]{\exists x\in X\,\exists y\in Y\,
    \exists z\in Z\;w=r(x,y,z)}{\rEE{8,9,14}}
  \proofstepwidestar[1,3]{\exists x\in X\,\exists y\in Y\,
    \exists z\in Z\;w=r(x,y,z)}{\rEE{5,6,15}}
  \proofstepwidestar[1,2]{\exists x\in X\,\exists y\in Y\,
    \exists z\in Z\;w=r(x,y,z)}{\rEE{2,3,16}}
  \proofstepwidestar[1]{\begin{aligned}[t]
    &(\exists x\in X\,\exists y\in Y\,\exists z\in Z\;
      w=t(x,y,z))\rightarrow{}\\[-2pt]
    &(\exists x\in X\,\exists y\in Y\,\exists z\in Z\;
      w=r(x,y,z))
  \end{aligned}}{\rRI{2,17}}
  \closeproofpartwide

  \proofpartwide{Vom zweiten Term zum ersten}
  \setcounter{proofstepnr}{18}
  \proofstepwidestar[19]{\exists x\in X\,\exists y\in Y\,
    \exists z\in Z\;w=r(x,y,z)}{\rA}
  \proofstepwidestar[20]{x\in X\land\exists y\in Y\,
    \exists z\in Z\;w=r(x,y,z)}{\rA}
  \proofstepwidestar[20]{x\in X}{\rAEa{20}}
  \proofstepwidestar[20]{\exists y\in Y\,\exists z\in Z\;
    w=r(x,y,z)}{\rAEb{20}}
  \proofstepwidestar[23]{y\in Y\land\exists z\in Z\;
    w=r(x,y,z)}{\rA}
  \proofstepwidestar[23]{y\in Y}{\rAEa{23}}
  \proofstepwidestar[23]{\exists z\in Z\;w=r(x,y,z)}{\rAEb{23}}
  \proofstepwidestar[26]{z\in Z\land w=r(x,y,z)}{\rA}
  \proofstepwidestar[26]{z\in Z}{\rAEa{26}}
  \proofstepwidestar[26]{w=r(x,y,z)}{\rAEb{26}}
  \proofstepwidestar[1,20,23,26]{t(x,y,z)=r(x,y,z)}{%
    \rRE{27,\rUE{\rRE{24,\rUE{\rRE{21,\rUE{1}}}}}}}
  \proofstepwidestar[1,20,23,26]{r(x,y,z)=t(x,y,z)}{%
    \FormulaRefAuto{a=b\vdash b=a}[thm]{29}}
  \proofstepwidestar[1,20,23,26]{w=t(x,y,z)}{\rIE{30,28}}
  \proofstepwidestar[1,20,23,26]{\exists x\in X\,\exists y\in Y\,
    \exists z\in Z\;w=t(x,y,z)}{%
    \rEI{\rAI{21,\rEI{\rAI{24,\rEI{\rAI{27,31}}}}}}}
  \proofstepwidestar[1,20,23]{\exists x\in X\,\exists y\in Y\,
    \exists z\in Z\;w=t(x,y,z)}{\rEE{25,26,32}}
  \proofstepwidestar[1,20]{\exists x\in X\,\exists y\in Y\,
    \exists z\in Z\;w=t(x,y,z)}{\rEE{22,23,33}}
  \proofstepwidestar[1,19]{\exists x\in X\,\exists y\in Y\,
    \exists z\in Z\;w=t(x,y,z)}{\rEE{19,20,34}}
  \proofstepwidestar[1]{\begin{aligned}[t]
    &(\exists x\in X\,\exists y\in Y\,\exists z\in Z\;
      w=r(x,y,z))\rightarrow{}\\[-2pt]
    &(\exists x\in X\,\exists y\in Y\,\exists z\in Z\;
      w=t(x,y,z))
  \end{aligned}}{\rRI{19,35}}
  \proofstepwidestar[1]{\begin{aligned}[t]
    &(\exists x\in X\,\exists y\in Y\,\exists z\in Z\;
      w=t(x,y,z))\leftrightarrow{}\\[-2pt]
    &(\exists x\in X\,\exists y\in Y\,\exists z\in Z\;
      w=r(x,y,z))
  \end{aligned}}{\rLRI{18,36}}
  \closeproofpartwide
\end{tabproofsplitwide}

\FormulaThmDeltaK[Mengengleichheit aus punktweise gleichen Dreifachtermen]{%
  \begin{aligned}[t]
    &\forall x\in X\,\forall y\in Y\,\forall z\in Z\;
      t(x,y,z)=r(x,y,z)\dsep{}\\[-2pt]
    &\forall w\,(w\in M\leftrightarrow
      \exists x\in X\,\exists y\in Y\,\exists z\in Z\;
      w=t(x,y,z))\dsep{}\\[-2pt]
    &\forall w\,(w\in N\leftrightarrow
      \exists x\in X\,\exists y\in Y\,\exists z\in Z\;
      w=r(x,y,z))\vdash M=N
  \end{aligned}
}{TripleTermRepresentationEquality}{%
  \DeltaRow{Mengen}{X\dsep Y\dsep Z\dsep M\dsep N\dsep
    x\dsep y\dsep z\dsep w}
  \DeltaRow{Mengenterme}{t(x,y,z)\dsep r(x,y,z)}
}
\begin{tabproofwide}
  \proofstepwidestar[1]{\forall x\in X\,\forall y\in Y\,
    \forall z\in Z\;t(x,y,z)=r(x,y,z)}{\rA}
  \proofstepwidestar[2]{\begin{aligned}[t]
    &\forall w\,(w\in M\leftrightarrow{}\\[-2pt]
    &\quad\exists x\in X\,\exists y\in Y\,\exists z\in Z\;
      w=t(x,y,z))
  \end{aligned}}{\rA}
  \proofstepwidestar[3]{\begin{aligned}[t]
    &\forall w\,(w\in N\leftrightarrow{}\\[-2pt]
    &\quad\exists x\in X\,\exists y\in Y\,\exists z\in Z\;
      w=r(x,y,z))
  \end{aligned}}{\rA}
  \proofstepwidestar[1]{\begin{aligned}[t]
    &(\exists x\in X\,\exists y\in Y\,\exists z\in Z\;
      w=t(x,y,z))\leftrightarrow{}\\[-2pt]
    &(\exists x\in X\,\exists y\in Y\,\exists z\in Z\;
      w=r(x,y,z))
  \end{aligned}}{\FormulaRefAuto{TripleTermWitnessPointwiseEquality}[thm]{1}}
  \proofstepwidestar[1,2]{\begin{aligned}[t]
    &w\in M\leftrightarrow{}\\[-2pt]
    &\quad\exists x\in X\,\exists y\in Y\,\exists z\in Z\;
      w=r(x,y,z)
  \end{aligned}}{\rLRS{4,\rUE{2}}}
  \proofstepwidestar[1,2,3]{w\in M\leftrightarrow w\in N}{%
    \rLRS{\rUE{3},5}}
  \proofstepwidestar[1,2,3]{\forall w\,(w\in M\leftrightarrow w\in N)}{%
    \rUI{6}}
  \proofstepwidestar[1,2,3]{M=N}{%
    \FormulaRefAuto{\forall x\,(x\in A\leftrightarrow x\in B)
      \eqvdash A=B}[thm]{7}}
\end{tabproofwide}

Die Existenz der dargestellten Mengen ist in diesem Satz vorausgesetzt.
Insbesondere ist er unmittelbar auf bereits definierte Mengenprodukte
anwendbar: Nur die punktweise Gleichheit der beiden erzeugenden Terme
bleibt dann algebraisch nachzuweisen.
